\section{Results}
\label{sec:results}

All results are derived from the analytical sweep of Section~\ref{sec:methods}
at fixed operating conditions ($u_{\mathrm{air},\infty} = 5\,\mathrm{m\,s^{-1}}$,
$v_{\mathrm{i}} = 0.5\,\mathrm{m\,s^{-1}}$). Their combination under one set of
constraint values is treated separately as a case study in
Section~\ref{sec:case-study}.

\subsection{Bundle-scale conductance map}
\label{sec:results-ko}

Figure~\ref{fig:ko} shows the package-scaled conductance
$k_{\mathrm{o}}A_{\mathrm{o}}$ over the $(d_{\mathrm{o}},\,\tau)$ plane for
aluminum and PA, with the respective minimum-wall constraints.
Both maps share a common color scale, so the absolute conductance gap is
directly visible.
Smaller outer diameters both raise $k_{\mathrm{o}}$ and increase the available
outer surface area at fixed package size; the resistance distribution behind
this trend is quantified in the next subsection.

\begin{figure}[htbp]
\centering
\includegraphics[width=\textwidth]{20_bundle_kA_alu_pa_portrait}
\caption{Package-scaled heat-transfer conductance
  $k_{\mathrm{o}}A_{\mathrm{o}}$ for aluminum and PA over the
  $(d_{\mathrm{o}},\tau)$ design space. Line and symbol conventions as defined
  in Sect.~\ref{sec:plot-conventions}; the $d_{\mathrm{o}}$ axis is logarithmic.}
\label{fig:ko}
\end{figure}

\pagebreak

\subsection{Thermal resistance distribution}
\label{sec:results-shares}

Figure~\ref{fig:shares} maps the three fractional resistances $\phi_{\mathrm{i}}$, $\phi_{\mathrm{w}}$, and $\phi_{\mathrm{o}}$
for aluminum and PA tubes over the $(d_{\mathrm{o}},\,\tau)$ design space.

\begin{figure}[htbp]
\centering
\includegraphics[width=\textwidth]{12_shares_alu_pa_portrait.pdf}
\caption{Fractional thermal resistance shares $\phi_{\mathrm{i}}$,
  $\phi_{\mathrm{w}}$, and $\phi_{\mathrm{o}}$ for Al and PA tubes. Line and
  symbol conventions as defined in Sect.~\ref{sec:plot-conventions}.}
\label{fig:shares}
\end{figure}

\pagebreak
The fractional resistance shares are defined as
\[
  \phi_j = \frac{R_j}{R_{\mathrm{tot}}}, \qquad j \in \{\mathrm{i},\,\mathrm{w},\,\mathrm{o}\},
\]
with the total resistance $R_{\mathrm{tot}} = R_{\mathrm{i}} + R_{\mathrm{w}} + R_{\mathrm{o}}$ and
$\phi_{\mathrm{i}} + \phi_{\mathrm{w}} + \phi_{\mathrm{o}} = 1$ by construction.

The outer convective resistance $\phi_{\mathrm{o}}$ constitutes the dominant fraction across
the small-diameter minimum-wall region.
Even for PA, where $\lambda_{\mathrm{w}} = 0.25\,\mathrm{W\,m^{-1}\,K^{-1}}$ represents an
880-fold reduction relative to aluminum, the wall resistance share $\phi_{\mathrm{w}}$
is $2.1$--$5.5\,\%$ along the PA minimum-wall line for
$d_{\mathrm{o}}=0.3$--$2.0\,\mathrm{mm}$.
When the outer convective resistance dominates, reducing wall conductivity by two orders
of magnitude attenuates $k_{\mathrm{o}}$ in proportion to $\phi_{\mathrm{w}}$ -- the penalty shrinks
with the wall-resistance share.
The corresponding outer share is $\phi_{\mathrm{o}}=83.8$--$85.4\,\%$ in this range.
At larger wall-thickness ratios, $\phi_{\mathrm{w}}$ grows and the geometry becomes less
attractive for compact cooling applications.

The inner convective share $\phi_{\mathrm{i}}$ is secondary but still relevant under the
conservative EGL5050 VDI G1 assumption; along the same line it
is $\phi_{\mathrm{i}}=9.6$--$13.2\,\%$, while the outer-side resistance remains
the largest term.



\FloatBarrier
\subsection{Coolant flow limits}
\label{sec:coolant-flow-limits}

Two coolant-side hydraulic limits are applied to the fixed-$v_{\mathrm{i}}$
sweep.
In Fig.~\ref{fig:coolant_throughput}, the geometrically similar pitch
assumption in Eq.~\eqref{eq:coolant_throughput} is shown. 

\begin{figure}[htbp]
\centering
\includegraphics[width=80mm]{19_coolant_throughput_lmin}
\caption{Coolant throughput diagnostic at fixed tube-side velocity. Line and
  symbol conventions as defined in Sect.~\ref{sec:plot-conventions}; the
  $d_{\mathrm{o}}$ axis is logarithmic.}
\label{fig:coolant_throughput}
\end{figure}

It makes the total flow depend
on the open-area ratio $(d_{\mathrm{i}}/d_{\mathrm{o}})^2$ under this
scaling, rather than on $d_{\mathrm{o}}$ alone.
At the fixed tube-side velocity used here, the increasing tube count at smaller
diameter is offset by the smaller flow area per tube; the associated wall-shear
effect is evaluated in the separate pressure-loss check below.

\Needspace{15\baselineskip}
\FloatBarrier
\subsection{Coolant pressure-loss limits}
\label{sec:coolant-pressure-loss-limits}

The pressure-loss map in Fig.~\ref{fig:dp_i_fric} follows the laminar
Hagen--Poiseuille scaling in Eq.~\eqref{eq:dp_i_laminar} over the PA
minimum-wall range.
Along the PA minimum-wall line, the affected region is confined to the smallest diameters:
for $d_{\mathrm{o}}\gtrsim 0.31\,\mathrm{mm}$ the straight-tube friction
pressure stays below $0.5\,\mathrm{bar}$.

\begin{figure}[htbp]
\centering
\includegraphics[width=80mm]{18_tube_side_friction_pressure_drop}
\caption{Straight-tube coolant pressure loss at fixed tube-side velocity. Line
  and symbol conventions as defined in Sect.~\ref{sec:plot-conventions}; the
  $d_{\mathrm{o}}$ axis is logarithmic.}
\label{fig:dp_i_fric}
\end{figure}

\FloatBarrier
\subsection{Tolerance-adjusted pressure-integrity check}
\label{sec:results-burst}

Figure~\ref{fig:burst} shows the tolerance-adjusted Lam\'e burst-pressure maps
for aluminum and PA over the $(d_{\mathrm{o}},\,\tau)$ plane.
Because $\sigma_{\mathrm{lim},\mathrm{Al}} / \sigma_{\mathrm{lim},\mathrm{PA}} = 20$, the PA isolines
are shifted strongly toward thicker walls compared with aluminum at the same
internal pressure.
In the thermally attractive small-diameter range, however, parts of the PA
window remain above typical automotive cooling-circuit operating pressures.
This is only a diagnostic observation; the integrated design-boundary map uses
the stricter criterion
$p_{\mathrm{b},\mathrm{tol}}\ge 6\,\mathrm{bar}$.
The comparison between rows shows the sensitivity to absolute tolerance:
at the PA minimum wall, the normal route leaves
$t_{\mathrm{loc},\mathrm{min}}=5\,\mu\mathrm{m}$, whereas the high-precision route leaves
$20\,\mu\mathrm{m}$.
At fixed $t_{\mathrm{min}}$ and $\Delta t$, smaller $d_{\mathrm{o}}$ increases
$(t_{\mathrm{min}}-\Delta t)/d_{\mathrm{o}}$ and raises the Lam\'e geometry factor.
The check shifts attention from isolated tube-wall burst to
bundle-level mechanisms outside this idealized calculation.

\begin{figure}[htbp]
\centering
\includegraphics[width=\textwidth]{15_burst_tolerance_alu_pa_portrait}
\caption{Tolerance-adjusted Lam\'e burst pressure
  $p_{\mathrm{b},\mathrm{tol}}$ for aluminum and PA tubes.
  Columns compare materials; rows compare
  $\Delta t=0.020\,\mathrm{mm}$ and $\Delta t=0.005\,\mathrm{mm}$. Line and
  symbol conventions as defined in Sect.~\ref{sec:plot-conventions}; the
  $d_{\mathrm{o}}$ axis is logarithmic.}
\label{fig:burst}
\end{figure}

\FloatBarrier
\subsection{Capillary-rise process constraint}
\label{sec:results-capillary-rise}

Figure~\ref{fig:capillary_rise} maps the predicted resin rise for aluminum and
PA columns under effective acceleration factors $G=1$, 5, and 10.
At $G=1$, sub-millimeter clear gaps generate process-relevant resin rise;
because Eq.~\eqref{eq:capillary_rise} scales as $1/G$, increasing
acceleration to $G=5$--$10$ reduces the predicted rise by the same factor.
For the benchmark pitch ratio $S_{\mathrm{L}}/\!d_{\mathrm{o}}=2.00$, the
$h_{\mathrm{cap}}=2\,\mathrm{mm}$ line at $G=10$ is near
$d_{\mathrm{o}}\approx 0.31\,\mathrm{mm}$ for the aluminum estimate and
$d_{\mathrm{o}}\approx 0.27\,\mathrm{mm}$ for the PA\,/\,process proxy. The two
capillary limits are much closer than a PA material-constant
interpretation would suggest. Compact pitch geometries that are thermally
attractive still require a controlled potting route, such as centrifugal
potting, as an independent process constraint.

\begin{figure}[htbp]
\centering
\includegraphics[width=\textwidth]{14_capillary_rise_alu_pa_portrait}
\caption{Predicted capillary rise $h_{\mathrm{cap}}$ for aluminum and PA
  potting gaps under effective acceleration factors $G=1$, 5, and 10. Line and
  symbol conventions as defined in Sect.~\ref{sec:plot-conventions}; the
  $d_{\mathrm{o}}$ axis is logarithmic.}
\label{fig:capillary_rise}
\end{figure}

\FloatBarrier
